\section{Usage}

\subsection{Conversions}
In this section, we describe how to access the underlying component library abstractions from corresponding
Dyninst abstractions. The component libraries (SymtabAPI, InstructionAPI, ParseAPI, and PatchAPI) often
provide greater functionality and cleaner interfaces than Dyninst, and thus users may wish to use a mix
of abstractions. In general, users may access component library abstractions via a convert function, which
is overloaded and namespaced to give consistent behavior. The definitions of all component library
abstractions are located in the appropriate documentation.

\code {
PatchAPI::PatchMgrPtr PatchAPI::convert(BPatch_addressSpace *);

PatchAPI::PatchObject *PatchAPI::convert(BPatch_object *);

ParseAPI::CodeObject *ParseAPI::convert(BPatch_object *);

SymtabAPI::Symtab *SymtabAPI::convert(BPatch_object *);

SymtabAPI::Module *SymtabAPI::convert(BPatch_module *);

PatchAPI::PatchFunction *PatchAPI::convert(BPatch_function *);

ParseAPI::Function *ParseAPI::convert(BPatch_function *);

PatchAPI::PatchBlock *PatchAPI::convert(BPatch_basicBlock *);

ParseAPI::Block *ParseAPI::convert(BPatch_basicBlock *);

PatchAPI::PatchEdge *PatchAPI::convert(BPatch_edge *);

ParseAPI::Edge *ParseAPI::convert(BPatch_edge *);

PatchAPI::Point *PatchAPI::convert(BPatch_point *, BPatch_callWhen);

PatchAPI::SnippetPtr PatchAPI::convert(BPatch_snippet *);

SymtabAPI::Type *SymtabAPI::convert(BPatch_type *);


\subsection{Using the API}

In this section, we describe the steps needed to compile your mutator and mutatee programs and to run
them. First we give you an overview of the major steps and then we explain each one in detail.

To use Dyninst, you have to:

\begin{enumerate}

\item {
Build and install DyninstAP: DyninstAPI can be installed a package system such as Spack or can be
compiled from source. Our github webpage contains detailed instructions for installing Dyninst:
https://github.com/dyninst/dyninst.
}

\item {
Create a mutator program: You need to create a program that will modify some other
program. For an example, see the mutator shown in Appendix A.
}

\item {
Set up the mutatee: On some platforms, you need to link your application with
Dyninst's run time instrumentation library. [NOTE: This step is only needed in the
current release of the API. Future releases will eliminate this restriction.]
}

\item {
Run the mutator: The mutator will either create a new process or attach to an existing
one (depending on the whether createProcess or attachProcess is used).

\end{enumerate}


\subsubsection{Creating a mutator program}

The first step in using Dyninst is to create a mutator program. The mutator program specifies the
mutatee (either by naming an executable to start or by supplying a process ID for an existing
process). In addition, your mutator will include the calls to the API library to modify the mutatee.
For the rest of this section, we assume that the mutator is the sample program given in Appendix A.

The following fragment of a Makefile shows how to link your mutator program with the Dyninst
library on most platforms:

\code{

# DYNINST_INCLUDE and DYNINST_LIB should be set to locations
# where Dyninst header and library files were installed, respectively

retee.o: retee.c
  \$(CC) -c \$(CFLAGS) -I\$(DYNINST_INCLUDE) retee.c -std=c++11

retee: retee.o
  \$(CC) retee.o -L\$(DYNINST_LIB) -ldyninstAPI -o retee -std=c++11
}

On Linux, the options \code{-lelf} and \code{-ldw} may be required at the link step. You will also
need to make sure that the \code{LD_LIBRARY_PATH} environment variable includes the directory that
contains the Dyninst shared library.

Since Dyninst uses the c++11 standard, you will also need to enable this option for your compiler.
For GCC versions 4.3 and later, this is done by specifying \code{-std=c++0x}. For GCC versions 4.7 and
later, this is done by specifying \code{-std=c++11}. Some of these libraries, such as libdwarf and
libelf, may not be standard on various platforms. Check the README file in dyninst/dyninstAPI for
more information on where to find these libraries.

Under Windows NT, the mutator also needs to be linked with the dbghelp library, which is included
in the Microsoft Platform SDK. Below is a fragment from a Makefile for Windows NT:

\code{
# DYNINST_INCLUDE and DYNINST_LIB should be set to locations
# where Dyninst header and library files were installed, respectively

CC = cl

retee.obj: retee.c
  \$(CC) -c \$(CFLAGS) -I\$(DYNINST_INCLUDE)/h

retee.exe: retee.obj
  link -out:retee.exe retee.obj \$(DYNINST_LIB)\textbackslash{}libdyninstAPI.lib \textbackslash{}dbghelp.lib

}

\subsubsection{Setting Up the Application Program (mutatee)}

On most platforms, any additional code that your mutator might need to call in the mutatee (for example
files containing instrumentation functions that were too complex to write directly using the API) can
be put into a dynamically loaded shared library, which your mutator program can load into the mutatee
at runtime using the loadLibrary member function of BPatch_process.

To locate the runtime library that Dyninst needs to load into your program, an additional environment
variable must be set. The variable \code{DYNINSTAPI_RT_LIB} should be set to the full pathname of the
run time instrumentation library, which should be:

NOTE: DYNINST_LIB should be set to the location where Dyninst library files were installed

\code{
\$(DYNINST_LIB)/libdyninstAPI_RT.so (UNIX)

\%DYNINST_LIB/libdyninstAPI_RT.dll (Windows)
}

\subsubsection{Running the Mutator}

At this point, you should be ready to run your application program with your mutator. For example,
to start the sample program shown in Appendix A:

\code{retee foo <pid>}

